\documentclass[11pt,a4paper]{article}
\usepackage[T1]{fontenc}
\usepackage[utf8]{inputenc}
\usepackage[french]{babel}
\frenchsetup{SmallCapsFigTabCaptions=false}
% The English language files of this Babel installation are broken, so the
% French language is loaded for hyphenation only and every caption is renamed.
\addto{\captionsfrench}{%
  \renewcommand{\contentsname}{Contents}%
  \renewcommand{\abstractname}{Abstract}%
  \renewcommand{\figurename}{Figure}%
  \renewcommand{\tablename}{Table}%
  \renewcommand{\refname}{References}%
}
\usepackage{lmodern,microtype,amsmath}
\usepackage[margin=23mm]{geometry}
\usepackage{booktabs,array,tabularx,longtable}
\usepackage{xcolor,tikz,enumitem,float}
\usepackage{hyperref}
\usetikzlibrary{arrows.meta,positioning,fit,calc,shapes.geometric}
\definecolor{ink}{HTML}{1F242C}
\definecolor{muted}{HTML}{5C626C}
\definecolor{line}{HTML}{D9DEE5}
\definecolor{accent}{HTML}{AF0032}
\definecolor{rose}{HTML}{FDF0F4}
\definecolor{agent}{HTML}{EDF1F8}
\definecolor{agentline}{HTML}{5A6B8C}
\definecolor{determ}{HTML}{F2F2F0}
\definecolor{warn}{HTML}{8A5A00}
\hypersetup{colorlinks=true,linkcolor=black,citecolor=black,urlcolor=ink,
 pdftitle={LANCE: where the project stands},
 pdfauthor={LANCE project},pdflang={en-GB}}
\urlstyle{same}
\def\UrlBreaks{\do\/\do\-\do\_\do\.\do\:}
\newcolumntype{Y}{>{\raggedright\arraybackslash}X}
\DeclareRobustCommand{\code}[1]{{\small\path{#1}}}
\newcommand{\source}[2]{\href{../../#1}{\nolinkurl{#1}}. #2}
% Engineering detail lives in the margin, so the argument stays readable.
\DeclareRobustCommand{\marginref}[1]{\marginpar{\raggedright\scriptsize\textcolor{muted}{#1}}}
\newcommand{\revision}{2026-11-03 11:30 (Europe/Brussels, UTC+01:00)}
\setlength{\parskip}{3pt}
\setlength{\emergencystretch}{2em}
\setlist{nosep,leftmargin=5mm,topsep=3pt}
\clubpenalty=10000 \widowpenalty=10000
\setcounter{tocdepth}{1}
% A clear boundary between the state (which is the point) and the evidence.
\usepackage{titlesec}
\titleformat{\section}[block]
  {\normalfont\sffamily\bfseries\Large\color{ink}}
  {\colorbox{ink}{\makebox[1.1em][c]{\color{white}\thesection}}}
  {0.8em}{}[{\vspace{3pt}{\color{line}\titlerule[0.6pt]}}]
\titlespacing*{\section}{0pt}{18pt}{12pt}

\title{LANCE\\[4pt]
\large Where the project stands\\[2pt]
\normalsize Checkpoint 1 state report as of 3 November 2026}
\author{LANCE project\\\small Written by specialist strands, independently reviewed and synthesised}
\date{Last revision: \revision\\[3pt]
\small Code at \code{546b911} (30 October 2026); no source change since.\\
No new experiment is reported: the measurements cited are those already recorded.}

\begin{document}
\maketitle

\begin{abstract}
LANCE is a platform for auditing authorised connected-object networks with
language-model agents, inside a sandbox. A scenario is deployed, an agent audits
it, its evidence is archived, and what it found is scored against a ground truth.
Developed in the context of the UMONS work of the NATO SPS G6259 programme.

Its benchmark compares what one wants to compare, under the same conditions. This
checkpoint documents one case, an automation against a model.

The project counts a result only where no one-line rule settles the case and
reasoning over several observations does. Until such situations are written and
filtered, nothing else counts.

The position is unchanged since the previous checkpoint: the platform is built,
instrumented and documented, and its chain of evidence is sound.
\end{abstract}

\tableofcontents
\clearpage

% === WHERE THE PROJECT STANDS: the whole answer, before the evidence ===
\section{Where the project stands}
\label{sec:etat}

\subsection{The project in one paragraph}

LANCE lets language-model agents audit IoT networks. It builds the network it
wants to audit from a description, hands the scenario to an agent, archives everything the agent does, and scores what it found against a
ground truth fixed in advance. The point is not that a model should win: it is
that an audit carried out by an agent can be run, re-run and verified.



The benchmark compares what one wants to compare under the same conditions; this
checkpoint documents one case, an automation against a model. Comparing the two is
not possible yet: no trial separates them enough to tell them apart, and the arm
that holds everything constant but the interpretation does not exist
(Figure~\ref{fig:policies}). The platform is built, the comparison is half
instrumented, and the material to measure it has not been found.

% ---------------------------------------------------------------------------
\subsection{What an audit does, and where the model decides}
\label{sec:parcours}

An audit runs in six phases. Each phase produces one document that the next one
reads, and nothing else crosses between them. Figure~\ref{fig:audit} is the map to
keep in mind for the rest of the report: it shows the six phases, what each
produces, and the point that matters: \emph{who} decides at each step.

\begin{figure}[H]
\centering
\begin{tikzpicture}[x=1mm, y=1mm,
  phase/.style={draw=line, fill=white, line width=.55pt, rounded corners=2pt,
    text=ink, align=center, font=\sffamily\footnotesize, inner sep=3pt,
    minimum height=10mm, minimum width=17mm, text width=14mm},
  model/.style={phase, draw=agentline, fill=agent},
  fixed/.style={phase, fill=determ},
  flow/.style={draw=muted, line width=.6pt, -{Latex[length=1.4mm,width=.9mm]}}]
\node[fixed] (p1) at (-60,0) {\textbf{1. Graph}\\[-1pt]{\tiny\normalfont map}};
\node[model] (p2) at (-36,0) {\textbf{2. Recon}\\[-1pt]{\tiny\normalfont probe}};
\node[model] (p3) at (-12,0)  {\textbf{3. Guess}\\[-1pt]{\tiny\normalfont hypotheses}};
\node[model] (p4) at (12,0)   {\textbf{4. Test}\\[-1pt]{\tiny\normalfont verify}};
\node[model] (p5) at (36,0)  {\textbf{5. Attack}\\[-1pt]{\tiny\normalfont reach}};
\node[fixed] (p6) at (60,0)  {\textbf{6. Report}\\[-1pt]{\tiny\normalfont write}};
\draw[flow] (p1.east) -- (p2.west);
\draw[flow] (p2.east) -- (p3.west);
\draw[flow] (p3.east) -- (p4.west);
\draw[flow] (p4.east) -- (p5.west);
\draw[flow] (p5.east) -- (p6.west);
\node[font=\sffamily\tiny, text=muted]    at (-60,-9) {code};
\node[font=\sffamily\tiny, text=agentline] at (-36,-9) {model};
\node[font=\sffamily\tiny, text=agentline] at (-12,-9)  {model};
\node[font=\sffamily\tiny, text=agentline] at (12,-9)   {model};
\node[font=\sffamily\tiny, text=agentline] at (36,-9)  {model};
\node[font=\sffamily\tiny, text=muted]    at (60,-9)  {code};
\node[font=\sffamily\tiny, text=muted, anchor=east] at (-72,-4.5) {who\ decides};
\end{tikzpicture}
\caption{An audit in six phases, and where the model's freedom actually sits. The
model decides at phases 2, 3, 4 and 5: which probes to run, which weaknesses to
hypothesise, how to test each one, and how far to push an intrusion. What the code
decides is what becomes verifiable: phase 1 projects its evidence file from the
graph tools and refuses a submission that misses a device. When the model's
deliverable at phase 5 is missing or invalid, the code writes the assessment from
the tool ledger, and a completed phase never certifies that the objectives were
reached. Phase 6 assembles the
report from section sheets, with an agent that declares no tool and runs a single
turn. The score is produced elsewhere, by an offline evaluator that recounts the
log against an oracle fixed before the run.}
\label{fig:audit}
\end{figure}

\textbf{The log is the source, not the model's conclusion.} Every tool call is
archived with a reference, a refusal or an exception included. A finding that
cites a piece of evidence which does not appear in the log is rejected. This is
what makes a result attributable: the report states what happened, not what the
policy believed happened.

\textbf{The oracle precedes the run.} The expected answer is fixed before the
policy runs, in a file the policy never receives. A finding that contradicts the
oracle invalidates the trial instead of correcting it. This precedence is what
prevents the answer from being adjusted after the fact.

\textbf{The score is recomputed, never taken on trust.} The SSH, MQTT, API and V0
evaluators receive only the briefing, the oracle, the journal and the final
decision, and they refuse a decision that cites an observation absent from the
journal. The benchmark score of a pipeline run does not perform that check; the
confirmed metrics it would need are not emitted by the harnesses today.
\paragraph{Where this lives.}
Phase registry \code{src/agent/phases/registry.py}; governance boundary
\code{src/agent/core/executor.py}; evaluators under \code{src/benchmark/} and
\code{benchmarks/experiments/}.

% ---------------------------------------------------------------------------
\subsection{What runs in the sandbox}

LANCE deploys a scenario, hands it to an engine, archives what the engine does,
and scores the result against a ground truth fixed in advance. Three engines are
in question here: versioned procedures, the scanner followed by its rules, and a
model given the scanner's observations alone. All three receive the same scenario
and the same observations; what differs is what they are allowed to do with it. The
third does not exist yet. The benchmark scores the full pipeline only: the other two
run without a score, so a comparison across engines still needs an evaluator both of
them go through.

\begin{figure}[H]
\centering
\begin{tikzpicture}[x=1mm, y=1mm,
  known/.style={draw=muted, fill=determ, rounded corners=2pt, line width=.55pt,
    text=ink, align=center, font=\sffamily\footnotesize, inner sep=4pt,
    minimum height=13mm, text width=34mm},
  missing/.style={known, draw=accent, fill=rose, dashed, line width=.9pt},
  tiny/.style={font=\sffamily\tiny, text=muted},
  arrow/.style={-{Latex[length=1.8mm,width=1.1mm]}, draw=muted, line width=.7pt}]
\node[known] (scripted) at (-40,0) {\textbf{scripted}\\[-2pt]{\tiny versioned procedures}};
\node[known] (rules) at (0,0)      {\textbf{rules}\\[-2pt]{\tiny scanner, then rules}};
\node[missing] (llm) at (40,0)     {\textbf{llm\_on\_evidence}\\[-2pt]{\tiny the model on the same observations}};
\node[tiny, below=2mm of scripted] {runs, not scored};
\node[tiny, below=2mm of rules] {runs, not scored};
\node[tiny, below=2mm of llm, text=accent] {\textbf{does not exist}};
\node[tiny, above=2.5mm of rules, text=accent, text width=80mm, align=center]
  {all three receive the same scenario and the same observations};
\draw[arrow] (-40,13) -- (-40,7);
\draw[arrow] (0,13) -- (0,7);
\draw[arrow] (40,13) -- (40,7);
\end{tikzpicture}
\caption{The three policies and what is missing. The third would give the model the
observations the rules already receive and nothing else: no rule findings, no
ground truth, no expected count of flaws. It does not exist, and without it a
comparison has two arms and measures verification rather than interpretation.}
\label{fig:policies}
\end{figure}

\textbf{What is missing, and why it matters.} An exploration where the model roams
a network freely changes too much at once: the targets chosen, the order of the
probes and the interpretation all differ, so a better score cannot be attributed
to reasoning. Giving the model the scanner's observations isolates one
variable: \emph{who interprets}. That is the question the programme asked, and the
third policy is the only one that answers it.

% ---------------------------------------------------------------------------
% ---------------------------------------------------------------------------
\subsection{What the benchmark measures}
\label{sec:unites}

A comparison needs a unit of observation, or ten similar scenarios count as ten
proofs. Two units are enough, and the distinction between them is what keeps the
count honest.

\begin{itemize}[leftmargin=1.2em,itemsep=.35em,topsep=.4em]
\item \textbf{A scenario} is one deployment that gets audited: the devices, the
services, the injected flaw and the links between them. Four scenarios that differ
only by the kind of vulnerable machine are four pieces of matter, not four proofs.
\item \textbf{A trial} is what is submitted to the closure test: one scenario, or a
few that differ by a single axis. It counts as \emph{one} observation, whatever the
number of runs it contains.
\item \textbf{The number of runs does not enter the count.} The same scenario
audited four times by the same model remains one data point, not four;
\emph{series} stays a usage word for those runs, with no admission rule of its own.
\end{itemize}

Earlier reports said \emph{family}; the word is retired because it named the
grouping and never what the grouping had to prove.

\begin{figure}[H]
\centering
\begin{tikzpicture}[x=1mm,y=1mm,
  t/.style={font=\ttfamily\scriptsize, text=ink, anchor=west},
  n/.style={font=\sffamily\scriptsize, text=muted, anchor=west},
  a/.style={font=\sffamily\scriptsize\itshape, text=muted, anchor=west},
  gl/.style={draw=line, line width=.5pt}]
\node[t, font=\ttfamily\scriptsize\bfseries] (E) at (0,0) {trial/};
\node[n] at (42,0) {what is submitted to the closure test};
\draw[gl] (2.5,-1.2) -- (2.5,-17.5);
\draw[gl] (2.5,-4) -- (6,-4);
\draw[gl] (2.5,-8) -- (6,-8);
\draw[gl] (2.5,-12) -- (6,-12);
\draw[gl] (2.5,-17.5) -- (6,-17.5);
\node[t] at (7,-4)  {scenario 1.yaml};
\node[t] at (7,-8)  {scenario 2.yaml};
\node[t] at (7,-12) {scenario 3.yaml};
\node[t] at (7,-17.5) {scenario 4.yaml};
\node[n] at (42,-4)  {the vulnerable machine is a router};
\node[n] at (42,-8)  {\ldots{} is a sensor};
\node[n] at (42,-12) {\ldots{} is a public API};
\node[n] at (42,-17.5) {\ldots{} a private one};
\node[a] at (7,-22) {these four form a series: same topology, one axis varies};
\node[a] at (7,-25.5) {and they count as \textbf{one} observation, not four};
\end{tikzpicture}
\caption{What a trial contains. A \textbf{scenario} is one deployment that gets
audited. Scenarios that differ by a single axis form a \emph{series}: here, only the
kind of vulnerable machine changes. A \textbf{trial} is what is submitted to the
closure test, and it counts as one observation whatever the number of scenarios and
runs it holds.}
\label{fig:units}
\end{figure}

\textbf{What the closure test is.} Before building a trial, one asks whether a rule
that solves the case can be written in a single line. If it can, the trial measures
the omission of a declaration rather than a limit of interpretation, and it is not
built. The test is cheap; running it before writing eight fixtures is the change of
method this month produced.

The choices that follow from this state are stated in
\code{objectifs/ETAT-CURANT.md}; none of them is taken yet.

% ---------------------------------------------------------------------------
\subsection{The state in one table}
\label{sec:table-etat}

\begin{table}[H]
\centering
\caption{The state of the project on 3 November 2026, and where the proof is.}
\label{tab:etat-chiffre}
\small
\begin{tabularx}{\linewidth}{@{}p{3.9cm}Yp{2.6cm}@{}}
\toprule
Question & Answer & \textbf{Proof} \\
\midrule
Is an audit traceable? & Yes. Every tool call is archived with a reference,
bounded, and a piece of evidence that cannot be archived stops the run instead of
leaving a silent gap. The guards are application-level. & table~\ref{tab:composants} \\
Does the harness run outside the laboratory? & Yes, without broker and without a
model, under a guard whose blind spots are declared. It instruments a contrast; it
measures nothing. & table~\ref{tab:composants} \\
Are the protocols' contracts sound? & Yes. Briefing, private oracle and offline
evaluator are separated, and the oracle is fixed before the run. & table~\ref{tab:composants} \\
Can results be delivered reproducibly? & Yes. Push, validate, deploy to one of
three instances; no execution on a workstation. Single points of failure declared.
& table~\ref{tab:composants} \\
Does the model beat the rules? & \textbf{Unknown, and not yet measurable.} & Section~\ref{sec:construit} \\
How much of the corpus needs interpretation? & Two counts from two artefacts,
never to be cited one for the other. & Section~\ref{sec:construit} \\
Is the measurement statistically valid? & The tool is repaired and refuses an
indefensible figure; the trials that would give it data do not exist. & Section~\ref{sec:construit} \\
\bottomrule
\end{tabularx}
\end{table}

{\footnotesize\color{muted}Proof for the first four rows: Section~\ref{sec:construit},
table~\ref{tab:composants}. The last three: the same section.}

\subsection{How far the project goes}
\label{sec:jusquou}

Each open point below is paired with the observable test that would show it
settled, which is the one thing a table of states cannot carry. Three of the four
depend on the first. A comparable arm, a written reference
automation and a corpus built for the question are worth building only once a
situation exists that a one-line rule does not settle; before that, each of them
measures the absence of that situation. \textbf{The first line is the only work that
can start today.}

This is where the report stops. Everything up to a verified measurement is built;
the thing to be measured is not. The next result the project can legitimately
produce is \emph{a trial that survives its closure test}, not a better score. It is also the only line of the table whose proof needs no new agreement on method.

\clearpage


% === WHAT IS BUILT, AND WHAT IS MEASURED ===
\section{What is built, and what is measured}
\label{sec:construit}

\textbf{In brief.} Nine components exist, each with a stated guarantee and a
stated limit. None of them proves that a model does better than a script, which
is the object of the programme. The detail is in \code{docs/} for the behaviour
and in the archived reports for the implementation; what follows is what a
committee needs in order to know whether the instruments can be trusted.

\begin{table}[tbp]
\centering
\caption{What exists, what it guarantees, and what it does not.}
\label{tab:composants}
\small
\begin{tabularx}{\linewidth}{@{}>{\raggedright\arraybackslash}p{3.2cm}Y>{\raggedright\arraybackslash}p{3.5cm}@{}}
\toprule
Component & Guarantee & Limit \\
\midrule
Execution boundary & every tool call is archived with a reference, bounded by
budgets, refused outside the target network & application-level guards, not a
system sandbox \\
Laboratory & a scenario deployed, configured and reset, then destroyed: a trial
starts from the same state & one host, one container, one runner \\
Delivery & push, validate, deploy to one of three instances; no execution on a
workstation & no automatic rollback after a failed deployment \\
Audit pipeline & six phases, one document per phase, the model decides where
interpretation is needed (Figure~\ref{fig:audit}) & two scenarios on one
laboratory are not mutually excluded \\
Evidence contracts & briefing, private oracle and offline evaluator separated;
the oracle is fixed before the run & four protocols only \\
Harnesses & one contract, one deterministic engine, one model engine, one evaluator per
protocol & the H2 engines are offline; other controllers need the laboratory \\
Sandbox & the deterministic engines run without broker, model or laboratory, under a
guard whose blind spots are declared & it instruments a contrast, it measures
nothing \\
Corpus & 288 ground-truth entries over the catalogue & built around what signature
rules find (table~\ref{tab:corpus2}) \\
Inference tool & one-sided p-value against the pre-registered threshold, no
figure below three trials & it has nothing to compute on \\
\bottomrule
\end{tabularx}
\end{table}

\subsection{What is measured, and what the measures say}

\textbf{The fidelity gate passes, partially.} The V0 control, a task solvable
without inference, succeeds on both engines three times out of three, with a
saturating control: it is readable for a repeated defect, not for a rare one. The
threshold frozen at 0.9 is computed separately by the gate and conditions nothing.

\textbf{One trial was measured, and its gain was absorbed.} H2, on 22 October
2026: eight fixtures, 25 messages, two readable runs. The first ran the two
engines of the time; the reference run that replaced it rotates three engines over
three repetitions, which gives 72 runs. On the raw variants the model scores 0.5,
the declared parser 0.0, and a third engine reading by convention covers 8 of 8:
it absorbs the gain
without declaring anything.

\textbf{One trial was built and rejected in an hour.} The vendor-register probe
of 30 October resisted six recorded reviews and fell to a rule that reads the
effect of a write on the neighbouring registers. The filter that would have
caught it costs minutes: if the rule settling the case can be written in one
line, a script will have it.

\textbf{The inference is repaired and idle.} The tool now computes a one-sided
p-value against the threshold, refuses to emit a figure below three trials and
requires five to decide. No set of five independent trials exists, so the
arithmetic has nothing to work on.

\textbf{The exploratory pilot published three pairs, not twelve.} The
reconciliation of 16 October plans twelve pairs. Eight attempts across seven pairs
were executed under a freeze: three were published as comparable and five were
incomparable. The campaign is not admissible.

\subsection{The corpus, and why it cannot answer the question}

\begin{table}[H]
\centering
\caption{Two counts, from two artefacts, never to be cited one for the other.}
\label{tab:corpus2}
\small
\begin{tabular}{@{}>{\raggedright\arraybackslash}p{3.4cm}>{\raggedright\arraybackslash}p{3.4cm}>{\raggedright\arraybackslash}p{3.4cm}>{\raggedright\arraybackslash}p{2.4cm}@{}}
\toprule
Reading & Decidable by a rule & Not decidable & Total \\
\midrule
Mechanical (binary \code{type} of each entry) & 282 by signature & 6 by logic (2.1\,\%) & 288 \\
Interpretive (category judged as a whole) & 243 (84.4\,\%) & 35 mixed, 4 chaining, 6 logic & 288 \\
\bottomrule
\end{tabular}
\end{table}

Both are recomputed from \code{benchmarks/labels/derived\_index.json} for the
mechanical reading and from the coverage audit of 12 October 2026 for the
interpretive one, which totals categories rather than entries. A duel on this
corpus would rank rules against rules, which is exactly what H2 observed.


% === WHAT REMAINS ===
\section{What remains}
\label{sec:reste}

\subsection{What the project can assert, and what it cannot}

An audit produces verifiable
traces: every tool call is archived and referenced, and the run stops when its
evidence log cannot be written. The guards, however, are application-level. A harness
can run outside the laboratory, without broker or model, under a guard that
declares its blind spots. It instruments a contrast; it measures nothing. A third
script arm absorbs a gain attributed to the model on the most documented trial but that holds for eight fixtures and one property. The corpus is mostly
detectable by signature, which makes a null difference expected by construction but that concerns this corpus, not the situations the programme targets.

Must not be written as established: the superiority of the model; a significant
gap, since no independent trial exists to compute one; generality, H2 covering
one schema, one question and one model; a cost gain; the validity of the V0
threshold, which the gate does not apply; and the integrity of the corpus, the
official healthy control being absent from the catalogue.

\subsection{What would settle each open point}
\label{sec:jusquou}

Each open point is paired below with the observable test that would show it
settled, which is the one thing a table of states cannot carry. Three of the four
depend on the first: a comparable arm, a written reference automation and a corpus
built for the question are worth building only once a situation exists that a
one-line rule does not settle. Before that, each of them measures the absence of
that situation. \textbf{The first line is the only work that can start today.}

\begin{table}[H]
\centering
\caption{Each blockage, what would cross it, and how one would know. No term here
needs the glossary: the table is meant to be read before it.}
\label{tab:jusquou}
\small
\begin{tabularx}{\linewidth}{@{}p{3.4cm}Yp{4.0cm}@{}}
\toprule
Blockage & What would cross it & How one would know \\
\midrule
No situation a rule cannot decide & a case where no one-line rule settles the answer &
the closure script runs and refuses to admit it \\
No comparable arm & the model interpreting the very observations the rules
receive & the same scenario scored on both, by the same evaluator \\
The competent reference is a doctrine & a written reference automation standing in
for it & it passes a trial that the model also passes \\
Corpus built for rules & scenarios where a well-written rule can only abstain &
the interpretive share rises, recomputed entry by entry \\
\bottomrule
\end{tabularx}
\end{table}


This is where the report stops. Everything up to a verified measurement is built;
the thing to be measured is not. The next result the project can legitimately
produce is \emph{a trial that survives its closure test}, not a better score. It
is also the only line of the table whose proof needs no new agreement on method.

\begin{samepage}
\section{Conclusion}

The system is built, instrumented and documented. Its chain of evidence is sound:
a separate oracle, an evaluator that recounts, a run that stops when its evidence
cannot be archived, a safeguard that declares its blind spots. The verification discipline works: every claim here was checked against the code,
not against what it says about itself.

The scientific question remains open. The programme asked when a model does better
than a script with identical tools; the available data do not answer it, and the
current arrangement would not allow an answer even with more machine time. The
next result the project can legitimately produce is \emph{a trial that survives its
closure test}, not a better score.
\end{samepage}


\end{document}
